% Abstract: <= 200 words (CJFR). Findings sentences are placeholders until the re-run is analysed.
\section*{Abstract}
Rolling-horizon optimization is widely used to schedule forest harvesting machines, but there is
little evidence on how long the planning horizon should be and how often the plan should be
re-optimized. We evaluated both choices with the open-source Forest Harvesting Operations Planning
System (FHOPS) on nine semi-synthetic British Columbia case studies (three timber supply areas,
6--40 blocks, 6--24 machines, 112-day horizon). The design crossed four planning horizons
(2--16 weeks), three re-optimization intervals (1, 7, and 14 days), and two solution methods
(simulated annealing and mixed-integer programming). Each window started from the operational
state left by the previous plan: remaining volume, staged inter-phase inventory, and machine
positions. Every run was scored by replaying its stitched, implemented schedule against the full
problem. \pending{Main findings: effect of re-optimization interval and horizon length on
delivered volume relative to full-horizon plans, sequencing feasibility, and computing time;
role of solver choice and tractability.} \pending{Implication for configuring operational
planning systems.} Evaluating rolling-horizon plans on their stitched schedules, with state
carried between windows, is essential: scoring single windows can misstate plan quality.

\vspace{1em}
\noindent\textbf{Key words:} forest operations planning; machine scheduling; rolling horizon;
mixed-integer programming; simulated annealing; open-source decision support

\clearpage
